%%%%%%%%%%%%%%%%%%%%%%%%%%%%%%%%%%
%%   SECTION 4: Explicit formula for the characteristic map      %%%
%%%%%%%%%%%%%%%%%%%%%%%%%%%%%%%%%%

\section{Explicit formulas for the characteristic map} 
\label{sec:characteristic}
%%%%%%%%%%%%%%%%%%%%%%%%%%%%%%%%%%
%%%%%%%%%%%%%%%%%%%%%%%%%%%%%%%%%%
%%%%%%%%%%%%%%%%%%%%%%%%%%%%%%%%%%

In this section we give an explicit diffeomorphism from the Grassmann manifold $\gras$ of oriented $2$-planes through the origin in $4$-space to a product of two $2$-spheres, and then use it to give a formula for the characteristic map $g_L:\tor\to S^2$ of a three-component link $L$ in the $3$-sphere. 

We also describe an alternative form $h_L:\tor\to S^2$ of the characteristic map, homotopic to $g_L$ but more convenient for the proof of  \thmref{A} that we will give in \secref{A}.   The formula for $h_L$ will reveal a close connection between the characteristic map and the classical Gauss map $T^2\to S^2$ for two-component links in $3$-space. 

Throughout we regard $\br^4$ as the algebra of {\itb quaternions}, with orthonormal basis $1,i,j,k$ and unit sphere $S^3$ (oriented so that $i,j,k$ is a positive frame for the tangent space to $S^3$ at the point $1$),  and $\br^3$ as the subspace of {\itb pure imaginary quaternions} spanned by $i$, $j$ and $k$, with unit sphere $S^2$.    Thus $S^3$ is viewed as the multiplicative group of unit quaternions, and $S^2$ as the subset of pure imaginary unit quaternions.  

For any quaternion $q = q_0+q_1i+q_2j+q_3k$, let $\bar q$ denote its conjugate $q_0-q_1i-q_2j-q_3k$, which coincides with $q^{-1}$ when $q\in S^3$, and let $\re(q) = q_0$ and $\Im(q)=q_1i+q_2j+q_3k$ denote its real and imaginary parts.  


 
%%%%%
\part{The Grassmann manifold $\gras$}
%%%%%

It is well known
% dating back to the work of Julius Pl\"ucker in the nineteenth century
that $\gras$ can be identified with a product of two $2$-spheres.  In particular, we will use the diffeomorphism
$$
\pi: \gras \longrightarrow S^2\times S^2  
$$
that maps the oriented plane $\pb ab$, spanned by an orthonormal $2$-frame $(a,b)$ in $\br^4$, to the point $(b\bar a,\bar ab)$ in $S^2\times S^2$.  Note that both coordinates 
$$
\pi_+\pb ab \ = \ b\bar a \qquad\text{and}\qquad \pi_-\pb ab \ = \ \bar ab
$$
do in fact lie in $S^2$ since right and left multiplication by $\bar a$ are orthogonal transformations of $\br^4$, carrying the orthonormal frame $(a,b)$ to the orthonormal frames $(1,b\bar a)$ and $(1,\bar ab)$.   

To see that $\pi$ is well-defined, consider any other orthonormal basis for the plane $\pb ab$.  It must be of the form $(ac,bc)$ for some $c$ in the circle subgroup $C_{\bar ab}$ through $\bar ab$, since this group acts on the plane by rotations.   Thus it suffices to check that $bc\,\overline{ac} = b\bar a$, which is immediate, and that $\overline{ac}\, bc = \bar ab$, which is true since $c$ commutes with $\bar ab$.  

To see that $\pi$ is in fact a diffeomorphism, we can simply write down the inverse.  It maps a pair $(x,y)$ in $S^2\times S^2$ to the plane $\pb c{cy}$ where $c$ is the midpoint of any geodesic arc from $x$ to $y$ on $S^2$.   This can be verified by a straightforward calculation using the fact that conjugation by a pure imaginary quaternion rotates the $2$-sphere about that quaternion by $\pi$ radians, so $cy\bar c = x$.
% together with the observation that $C_{\bar ab}$ is the full centralizer of $\bar ab$ in $S^3$

\begin{comment}
More generally, conjugation by $\cos\theta + u\sin\theta$ rotates $S^2$ about $u$ by $2\theta$ radians.

Using this special property of the midpoint $c$, it is immediate that the map above, call it $\rho$, is a right inverse to $\pi$, since $\pi\rho(x,y) = (cy\bar c, \bar ccy) = (x,y)$.  To see that it is also a left inverse, we compute $\rho\pi\pb ab = \pb c{c\bar ab}$ where $c$ is the midpoint of any geodesic from $b\bar a$ to $\bar ab$ on $S^2$.  But then $c\bar ab\bar c = b\bar a$ by the special property of $c$, and it follows that $\bar ac$ commutes with $\bar ab$.  Therefore $c\in aC(\bar ab)$ and so $\rho\pi\pb ab \ = \  \pb c{c\bar ab} \ = \  \pb ab$.
\end{comment}

%%%%%%%%%%%%%%%%%%%%%%%%
%%%%%%%%%%%%%%%%%%%%%%%%
\part{The characteristic map $g_L$} 
%%%%%%%%%%%%%%%%%%%%%%%%
%%%%%%%%%%%%%%%%%%%%%%%%

Recall from the introduction that the Grassmann map $G:\conf\to\gras$ sends a triple $(x,y,z)$ of distinct points in $S^3$ to the plane they span in $\br^4$, translated to pass through the origin, and oriented so that
$$
G(x,y,z) \ = \ \pb{x-z}{y-z}.
$$  
We have extended notation so that for {\it any} two linearly independent vectors $a$ and $b$ in $\br^4$, the symbol $\pb ab$ denotes the oriented plane they span.  Then, given a three-component link $L$ in $S^3$, its characteristic map $g_L:\tor\to S^2$ is defined using the Grassmann map $G$ and the projection $\pi_+:\gras\to S^2$ by the formula 
$$
g_L(s,t,u) \ = \ \pi_+G(x,y,z) \ = \ \pi_+\pb{x-z}{y-z}
$$
where $x=x(s)$, $y=y(t)$ and $z=z(u)$ parametrize the components of $L$.   

To make this explicit, and to show that using $\pi_-$ in place of $\pi_+$ in the definition would not change the homotopy class of $g_L$, we need expressions for $\pi_\pm\pb ab$ when $a$ and $b$ are arbitrary linearly independent vectors in $\br^4$, but not necessarily orthonormal.  

For example, if $a$ and $b$ are orthogonal, then $b\bar a$ and $\bar ab$ are still pure imaginary, and so need only be normalized to give  $\pi_+\pb ab$ and $\pi_-\pb ab$.   

For a general pair of linearly independent vectors $a$ and $b$, the vector $c = b-(b\dt a/a\dt a)a$, where $\dt$ is the dot product in $\br^4$, is orthogonal to $a$ and satisfies $\pb ac = \pb ab$.  Therefore $\pi_+\pb ab = \pi_+\pb ac$, which equals the unit normalization of the vector $c\bar a = b\bar a - b\dt a$.  But $b\dt a = \re(b\bar a)$, and so $\pi_+\pb ab$ is the unit normalization of $\Im(b\bar a)$.  Similarly $\pi_-\pb ab$ is the unit normalization of $\Im(\bar ab)$.  Therefore, for {\it any} two linearly independent vectors $a$ and $b$, we have  
$$
\pi_\pm\pb ab \ = \ (a,b)_\pm / |(a,b)_\pm|  
$$
where $(\,,\,)_\pm$ are the skew symmetric bilinear forms on $\br^4$ defined by 
$$
(a,b)_+ \ = \ \Im\,b\bar a \qquad\text{and}\qquad (a,b)_- \ = \ \Im\,\bar ab\,.
$$

\begin{comment}
Note that both vectors $(a,b)_\pm$ have norm equal to the area of the parallelogram spanned by $a$ and $b$.
\end{comment}

It follows that $g_L(s,t,u)$ is the unit normalization of the vector
$$
\begin{aligned}
F(x,y,z) \ &= \ (x-z,y-z)_+ \ = \ (x,y)_+ + (y,z)_+ + (z,x)_+ \\
&= \ (ix\dt y + iy\dt z + iz\dt x \ , \  jx\dt y + jy\dt z + jz\dt x \ , \  kx\dt y + ky\dt z + kz\dt x)
\end{aligned}
$$
where the last expression follows from the formula $(a,b)_+ = (ia\dt b, ja\dt b, ka\dt b)$.  This formula is seen as follows.  By definition, $(a,b)_+ = \Im(b\bar a)  =  (i\dt b\bar a)i + (j\dt b\bar a)j + (k\dt b\bar a)k$.  But $q\dt b\bar a = qa\dt b$ for any $q$ (in particular $i$, $j$ or $k$) since right multiplication by the unit quaternion $a/|a|$ is an isometry, and so $(a,b)_+ = (ia\dt b)i + (ja\dt b)j + (ka\dt b)k = (ia\dt b, ja\dt b, ka\dt b)$.
% Similarly $\Im(\bar a b)= (ai\dt b, aj\dt b, ak\dt b)$
Summarizing, we have shown:
 
%%%%%%%%%
\begin{proposition}
\label{prop:char}
\textbf{\boldmath The characteristic map $g_L:\tor\to S^2$ of a three-component link $L$ in $S^3$ is given by the formula
$$
g_L(s,t,u) \ = \ F(x,y,z)/\|F(x,y,z)\|
$$
where $x = x(s)$, $y = y(t)$ and $z = z(u)$ parametrize the components of $L$, and $F:\conf\to\br^3$ is the function defined above.}
\end{proposition}
%%%%%%%%%

This formula for $g_L$ will be used in our proof of \thmref{B} in Sections \ref{sec:B} and \ref{sec:B'}.  
For \thmref{A} it will be more convenient, for the most part, to use an alternative form of the characteristic map that we introduce next.

%%%%%%%%%%%%%%%%%%%%%%%%
%%%%%%%%%%%%%%%%%%%%%%%%
\part{The asymmetric characteristic map $h_L$}
%%%%%%%%%%%%%%%%%%%%%%%%
%%%%%%%%%%%%%%%%%%%%%%%%

Given a three-component link $L$ in the $3$-sphere, we define below an asymmetric version 
$$
h_L:\tor\longrightarrow S^2
$$
of the characteristic map in which the last component of $L$ plays a distinguished role, and show that it is homotopic to the earlier defined symmetric characteristic map $g_L$.  As will be seen, the map $h_L$ can be viewed as a parametrized family of Gauss maps for twisted versions of the first two components of $L$, parametrized by the third component.  

For notational economy, we will use $\norm{q}$ to denote the unit normalization $q/|q|$ of any nonzero quaternion $q$.

The key observation that motivates the definition of $h_L$ is that the Grassmann map $G:\conf\to\gras$ factors up to homotopy through the {\itb Stiefel manifold} $\stiefel$ of orthonormal $2$-frames in $4$-space.  

More precisely, let $\pr_z$ denote {\itb stereographic projection} of $S^3-\{z\}$ onto $z^\perp$, the $3$-plane through the origin in $\br^4$ orthogonal to $z$.  Then we define the {\itb Stiefel map}
$$
H:\conf\longrightarrow\stiefel \quad,\quad (x,y,z) \longmapsto (z,\norm{\pr_zx-\pr_zy})
$$
(recall that the square brackets signify unit normalization), and will show that it is a homotopy equivalence whose composition with the canonical projection 
$$
P:\stiefel\longrightarrow\gras \quad,\quad (z,v)\longmapsto\pb zv
$$
is homotopic to the Grassmann map $G$.  

To see this, first observe that there is a deformation retraction of $\conf$ to its subspace 
$$
V  \ = \ \{(v,-z,z) \st v\perp z\}
$$ 
defined as follows.  Start with $(x,y,z)\in\conf$, and consider the points $\pr_zx$ and $\pr_zy$ in $z^\perp$.   Translation in $z^\perp$ moves $\pr_zy$ to the origin, and then dilation in $z^\perp$ makes the translated $\pr_zx$ into a unit vector.   Conjugating this motion by $\pr_z$ moves $x$ to $\norm{\pr_zx-\pr_zy}$, moves $y$ to $-z$, and leaves $z$ fixed, as pictured in \figref{stfl}, thus defining a deformation retraction of $\conf$ to its subspace $V$, sending $(x,y,z)$ to $(\norm{\pr_zx-\pr_zy},-z,z)$.
  

%%% FIGURE 14: Deformation retraction to V %%%
\begin{figure}[h!]
\includegraphics[height=150pt]{FigRT.pdf}
\put(-127,136){\small$z$}
\put(-105,18){\small$-z$}
\put(-80,103){\small$x$}
\put(-92,123){\small$y$}
\put(-40,72){\small $\pr_zx$}
\put(-30,100){\small $\pr_zy$}
\put(-187,45){$z^\perp$}
\put(-145,90){$S^3$}
\put(-150,0){\small$\norm{\pr_zx-\pr_zy} \ = \ (\pr_zx-\pr_zy)/|\pr_zx-\pr_zy|$}
\caption{The deformation retraction $\conf\to V$}
\label{fig:stfl}
\end{figure}
%%%%%%%%%%%%%%%%% 

Identifying $V$ with $\stiefel$ via $(v,-z,z)\leftrightarrow(z,v)$, this shows that $H$ is a homotopy equivalence with homotopy inverse $I$ given by $I(z,v) = (v,-z,z)$.  The calculation 
$$
GI(z,v) \ = \ \pb{v-z}{-2z} \ = \ \pb zv \ = \ P(z,v)\,,
$$
shows that $GI=P$, and so $G\simeq PH$ as claimed.

Now observe that there are two natural homeomorphisms from the Stiefel manifold $\stiefel$ to $S^3\times S^2$ that arise from viewing $S^3$ as the unit quaternions and $S^2$ as the pure imaginary unit quaternions, namely $(z,v) \mapsto (z,v\bar z)$ and $(z,v)\mapsto (z,\bar zv)$.  These yield projections $\pi_\pm:\stiefel\to S^2$ given by 
$$
\pi_+(z,v) \ = \ v\bar z \qquad\text{and}\qquad \pi_-(z,v) =  \bar zv.
$$
These are just the lifts to $\stiefel$ of the previously defined projections $\pi_\pm:\gras\to S^2$ with the same names. 
 
Define the {\itb asymmetric characteristic map} $h_L:\tor\to S^2$ to be the composition $\pi_+He_L$, that is,
$$
h_L(s,t,u) \ = \ \norm{\pr_zx-\pr_zy}\, \bar z 
$$ 
where as usual $x=x(s)$, $y=y(t)$ and $z=z(u)$ parametrize the components of $L$, and the square brackets indicate unit normalization.     

%%%%%%%%%
\begin{proposition}
\label{prop:asymmetric}
\textbf{\boldmath The two versions $g_L$ and $h_L$ of the characteristic map of a three-component link $L$ in $S^3$ are homotopic.  Furthermore, these maps are independent, up to homotopy, of the choice of projections $\pi_+$ or $\pi_-$ used in their definitions.}
\end{proposition}
%%%%%%%%%

\begin{proof}
By definition $g_L = \pi_+Ge_L$ and $h_L = \pi_+He_L$, shown in the diagram below as the maps from left to right across the bottom and top, respectively.
\vskip-.1in
\begin{diagram}[size=1.5em]
&&&& \quad \stiefel \quad && \\
&&& \ruTo[leftshortfall=10pt]^H \ && \rdTo[leftshortfall=10pt]^{\pi_+} & \\
\tor\ & \rTo[l>=.5in]_{e_L} & \ \conf \ \ \  && \dTo_P && \ S^2\\
&&& \rdTo[leftshortfall=10pt]_G && \ \ruTo[leftshortfall=7pt]_{\pi_+}[rightshortfall=10pt] & \\
&&&& \quad \gras \quad && \\
\end{diagram}
\vskip.1in
\noindent Here $e_L(s,t,u)=(x(s),y(t),z(u))$ records the parametrization of the link, and $G$ and $H$ are the Grassmann and Stiefel maps with their associated projections $\pi_+$.  It was shown above that the left triangle in the diagram commutes up to homotopy, and the right triangle commutes on the nose.  Therefore $h_L$ and $g_L$ are homotopic.

%\quad\underset{e_L}\longrightarrow\quad


Now the same argument shows that the maps $h_L' = \pi_-He_L$ and $g_L' =  \pi_-Ge_L$ are homotopic.   Hence to complete the proof, it suffices to show that $h_L$ and $h_L'$ are homotopic.  To do so, view $\stiefel$ as the unit tangent bundle of $S^3$, with projection $p$ to the base $S^3$ given by $p(z,v) = z$.  Then the composition $pHe_L$ is null-homotopic, since it maps onto the third component of $L$, and so the map $He_L$ is homotopic to a map into any $S^2$-fiber of the bundle $\stiefel\to S^3$.  We choose the fiber over $z=1$, where the two projections $\pi_+(z,v) = v\zbar$ and $\pi_-(z,v) = \zbar v$ coincide.  It follows that $h_L$ and $h_L'$ are homotopic. 
\end{proof}

%%%%%%%%%%%%%%%%%%%%%%
\part{A Gaussian view of the asymmetric characteristic map}
%%%%%%%%%%%%%%%%%%%%%%

The formula above for $h_L$ involves first normalizing a vector in $z^\perp$, and then multiplying by $\zbar$ to move it to the unit sphere $S^2$ in the pure imaginary quaternion $3$-space $\br^3$.  We would like to express this directly as the normalization of a vector in $\br^3$.  

A geometric argument shows that stereographic projection $\pr_a$ is given by 
$$
\pr_ab \ = \ (\Im\,b\bar a)a/(1-\re\,b\bar a)
$$
and it follows that $(\pr_ab)c  = \pr_{ac}\,bc$ for any three unit quaternions $a$, $b$ and $c$ with $a\ne b$.  Hence the formula $h_L(s,t,u) = \norm{\pr_zx-\pr_zy}\, \bar z$ can be rewritten as
$$
\begin{aligned}
h_L(s,t,u) \ &= \ \norm{\pr_1\,x\zbar - \pr_1\,y\zbar} \\
&= \ \norm{\pr_{-1}(-y\zbar) - \pr_{-1}(-x\zbar)}
\end{aligned}
$$
where, as usual, $x=x(s)$, $y=y(t)$ and $z=z(u)$ parametrize $L$, and the square brackets indicate unit normalization in $\br^3$.  We favor the last expression because stereographic projection from $-1$ is orientation-preserving, while from $1$ it is orientation-reversing.  

Now for any two distinct unit quaternions $a$ and $z$, introduce the abbreviation
$$
a_z \ = \ \pr_{-1}(-a\zbar) \ \in \ \br^3\,,
$$
and so in particular $a_{-1} = \pr_{-1}\,a$.  Then we can write
$$
h_L(s,t,u) \ = \ \norm{y_z-x_z}\,.
$$
For fixed $u$, this is just the classical Gauss map for the two-component link $X_z\cup Y_z \subset \br^3$ that is the image of the first two components $X\cup Y$ of $L$ under the map $a\mapsto a_z$.  
Thus the asymmetric characteristic map $h_L$ can be viewed as a one-parameter family of Gauss maps for images in $3$-space of $X\cup Y$.  The third component $Z$ provides the parameter and determines the axes about which these images are gradually twisted.

In the next section, we will explain this perspective more carefully.  But we can see right now that it yields an easy proof of the first part of \thmref{A}, equating the pairwise linking numbers of the link $L$ with the degrees of the restriction of its characteristic map $g_L$ to the coordinate $2$-tori.

%%%%%%%%%%%%%%%%%%%%%%
%%%%%%%%%%%%%%%%%%%%%%
\part{Proof of the first statement in \thmref{A}}
%%%%%%%%%%%%%%%%%%%%%%
%%%%%%%%%%%%%%%%%%%%%%

It can be arranged by an isotopy of $L$ that $z(0) = -1$.  Then the restriction of $h_L$ to the coordinate $2$-torus $S^1\times S^1\times 0$ is precisely the Gauss map of the stereographic image $X_{-1}\cup Y_{-1}$ of $X\cup Y$, whose degree is equal to the linking number $\lk(X, Y )$ since $\pr_{-1}$ is orientation-preserving. Since $g_L$ is homotopic to $h_L$, the same is true for $g_L$.  But then it follows from the symmetry of $g_L$ that $\lk(X,Z)$ and $\lk(Y,Z)$ are given by the degrees of $g_L$ on $S^1\times 0\times S^1$ and on $0 \times S^1\times S^1$, respectively. 

The proof of the second statement in \thmref{A} -- which relates the triple linking number of $L$ to the Pontryagin $\nu$-invariant of its characteristic map -- is more delicate.  It will occupy us for the next two sections.










\begin{comment}

%%%%%%%%%%%%%%%%%%%%%%
\part{A Gaussian view of the asymmetric characteristic map}
%%%%%%%%%%%%%%%%%%%%%%

Recall that the linking number of a two-component link in $\br^3$ is given by the degree of its Gauss map $(s,t)\mapsto \norm{y-x}$, where $x=x(s)$ and $y=y(t)$ parametrize the components of the link, and where as usual the brackets denote unit normalization.

Using stereographic projection $\pr_1:S^3-\{1\} \to \br^3$, the same formula holds for a two component links\ $L$ in $S^3-\{1\}$ where the associated {\itb Gauss map} 
$$
g_L:T^2\longrightarrow S^2 
$$ 
is given by $g_L(s,t) = \norm{\pr_1x - \pr_1y}$.  The change in the order of the variables $x$ and $y$ in the formula reflects the fact that $\pr_1$ is orientation reversing.  Here one uses the usual conventions for orienting $S^3\subset\br^4$ (outward normal first) and the pure quaternion $3$-space (with ordered basis $i,j,k$).

Now let $L$ be a $3$-component link in $S^3$, with components $X$, $Y$ and $Z$ parametrized by $x=x(s)$, $y=y(t)$ and $z=z(u)$ as usual.  For each $u\in S^1$, let $u(X\cup Y)$ denote the two-component link  that is the image of $X\cup Y$ under right multiplication by $\zbar = \zbar(u)$.  Then the asymmetric characteristic map of $L$ can be expressed as
$$
h_L(s,t,u) \ = \   \norm{\pr_1\,x\zbar - \pr_1\,y\zbar}\ = \ g_{u(X\cup Y)}(s,t)
$$
and so can be viewed as a parametrized family of Gauss maps for twisted versions $u(X\cup Y)$ of $X\cup Y$, where the twisting is induced by $Z$.

From this perspective, the first statement in \thmref{A} is completely transparent.  Indeed, it can be arranged by an isotopy of $L$ that $z(1) = 1$.  Then the restriction of $h_L$ to the coordinate $2$-torus $S^1\times S^1\times 1$ is precisely the Gauss map of $X\cup Y$, whose degree is equal to the linking number $\lk(X, Y )$. Since $g_L$ is homotopic to $h_L$, the same is true for $g_L$.  But then it follows from the symmetry of $g_L$ that $\lk(X,Z)$ and $\lk(Y,Z)$ are given by the degrees of $g_L$ on $S^1\times 1\times S^1$ and $1 \times S^1\times S^1$, respectively. 

The proof of the second statement in \thmref{A} is more delicate, and will occupy us for the next two sections.

\end{comment}  

